% !TeX program = lualatex
% Agregátor reálných otázek SZZ pro okruh I3. Jednotlivé otázky jsou v souborech q_*.tex.
% Sestavení: ./build.sh 02_spolecna_informatika/I3_programovaci_jazyky/priklady   (nebo cd .../priklady && latexmk)
\documentclass[11pt,a4paper]{article}
\usepackage{szzpriklady}
% Build běží z adresáře priklady/ (CWD), obrázky jsou v src/fig/.
\graphicspath{{src/fig/}{fig/}}

\title{Příklady otázek -- Programovací jazyky\\[0.3em]\large Reálné otázky ze SZZ 2022--2025 (I3)}
\author{Jakub Klesa}
\date{Pokryté předměty: NPRG013 (Java) / NPRG035 (C\#) / NPRG041 (C++) - dle volby jazyka}

\begin{document}
\maketitle
\tableofcontents
\bigskip

\begin{poznamka}
Tento dokument obsahuje reálné otázky z minulých termínů SZZ (2022--2025) zařazené do
okruhu I3. U otázek s obrázkem je grafika oříznutá z originálního PDF (viz odkaz na
zdroj). Text, číselná data a tabulky jsou přepsané věrně.
Zadání i oficiální náčrty řešení pocházejí ze sad zveřejněných na webu MFF UK:\par
{\raggedright\url{https://www.mff.cuni.cz/cs/studenti/bakalarske-studium/statni-zaverecne-zkousky/bakalarske-statni-zkousky-studijniho-programu-informatika}\par}
Náčrty označené jako doplněné jsou vlastní, oficiální sada je k~dané otázce neobsahuje;
vlastní řešení jsou v~C++ (u~otázky podzim 2022 v~Javě) a~jejich kód byl ověřen překladem.
Pojmy odpovídají přednáškám D.~Bednárka \emph{Programování v~C++} a~\emph{Pokročilé programování v~C++}
(\url{https://teaching.mff.cuni.cz/nprg041-web/}). Fialové rámečky \emph{Lehčí v~C\#} ukazují
kratší odpověď v~C\#, který zadání také dovolují.
\end{poznamka}
\bigskip

% === Otázky (chronologicky). Doplň \input řádky a soubory q_*.tex do src/. ===
% Příklad: \input{q_2024-leto-28_bayesovske-site}
% ZACATEK-OTAZEK
\input{src/q_2022-leto-02_preklad-konstrukci}
\input{src/q_2022-podzim-01_programovani-rozhrani-grafu}
\input{src/q_2023-jaro-02_objektovy-navrh-bitmapovy-editor}
\input{src/q_2023-leto-18_komponentovy-system}
\input{src/q_2024-podzim-02_reprezentace-typu}
\input{src/q_2025-jaro-02_knihovna-matice}
% KONEC-OTAZEK

\end{document}
